\section{Revue de la littérature}
\label{sec:revue}

Cette section passe en revue les travaux théoriques et empiriques pertinents
pour notre problématique. La revue théorique présente les théories des
transferts de migrants et les canaux par lesquels ils peuvent réduire les
privations infantiles. La revue empirique recense ensuite les résultats de la
littérature et dresse le bilan des lacunes que ce mémoire entend combler.

% ── PARTIE I : Revue théorique ──────────────────────────────
\subsection{Revue théorique}

Cette revue théorique établit les fondements de notre analyse. Elle présente
les définitions et l'ampleur des transferts de migrants, puis les théories qui
expliquent les raisons de la migration et les motivations à transférer des
fonds.

\subsubsection{Transferts de migrants : définition, ampleur et raisons de la migration}

Cette sous-section définit les transferts de migrants et en présente l'ampleur
mondiale et sénégalaise, avant d'exposer les principaux cadres théoriques
qui expliquent les motivations à transférer.

Les transferts de fonds des migrants, communément appelés \textit{remittances} dans
la littérature anglophone, désignent les fonds envoyés par des personnes résidant
à l'étranger vers leur pays ou leur région d'origine. Le Fonds Monétaire International
distingue trois composantes principales : les transferts courants des travailleurs
migrants, les rémunérations des salariés et les transferts de patrimoine des migrants.
Dans la présente étude, nous retenons une définition large incluant l'ensemble des
envois de fonds privés, formels et informels, destinés aux ménages bénéficiaires
résidant au Sénégal.

La mesure des transferts se heurte à d'importantes limites statistiques.
Les flux transitant par des canaux informels, réseaux de confiance, transport
physique d'espèces, échappent aux statistiques officielles. Au Sénégal, selon
\textcite{ocde2017}, les transferts informels pourraient représenter entre 30 et 50\%
du total des flux reçus. Cette sous-estimation systématique doit être gardée à l'esprit
lors de l'interprétation des résultats empiriques.


Les transferts de fonds constituent l'un des flux financiers les plus importants vers
les pays en développement, surpassant depuis 2010 les flux d'aide publique au
développement \parencite{maimbo2005}. À l'échelle mondiale, \textcite{worldbank2023} estimait ces flux à
près de 794 milliards de dollars en 2022, dont 626 milliards à destination des
pays à revenu faible et intermédiaire. Pour l'Afrique subsaharienne, ces flux
atteignent environ 53 milliards de dollars en 2022, représentant en moyenne
3 à 5\% du PIB régional, mais avec de fortes hétérogénéités nationales.

Le Sénégal figure parmi les pays d'Afrique de l'Ouest les plus dépendants de
ces flux~: les transferts y représentent structurellement 10 à 12\,\% du PIB,
dépassant les IDE et l'aide publique au développement réunis \parencite{worldbank2023}.
La résilience de ces flux a été démontrée lors de la crise sanitaire~: après
une contraction de 1,6\,\% en 2020, ils ont rebondi de 11,8\,\% en 2021
\parencite{worldbank2020}.


Avant d'examiner les motivations à transférer, il faut comprendre pourquoi
les individus migrent, car la décision de migration et celle d'envoyer des
fonds sont indissociables. La théorie néoclassique fournit le premier cadre
explicatif de la migration. Dans sa version macroéconomique, la migration
résulte des différentiels de salaires entre régions ou pays~: la main-d'\oe uvre
se déplace des zones où le travail est abondant et faiblement rémunéré vers
celles où il est rare et mieux payé, jusqu'à l'égalisation tendancielle des
rémunérations. Dans sa version microéconomique, formalisée par
\textcite{sjaastad1962}, la migration est un investissement en capital humain~:
l'individu rationnel migre si la valeur actualisée des gains attendus dans la
région de destination, nette des coûts monétaires et psychologiques du
déplacement, excède celle de ses gains dans la région d'origine.
\textcite{harris1970} enrichissent ce cadre en montrant que la décision se fonde
sur le revenu \emph{espéré} plutôt que sur le revenu effectif~: un salaire
urbain élevé peut attirer des migrants même lorsque la probabilité d'obtenir
un emploi est faible, ce qui explique la persistance de la migration vers des
marchés du travail saturés. Dans cette perspective, les envois de fonds sont
le rendement de l'investissement migratoire~: le migrant transfère une partie
du surplus de revenu obtenu grâce au différentiel de salaires, et les
transferts croissent avec la réussite économique du migrant à destination.
Cette lecture éclaire le cas sénégalais, où les destinations privilégiées
(France, Italie, Espagne, États-Unis) offrent des écarts de rémunération
considérables avec le marché du travail local.

La théorie des réseaux complète cette approche individuelle en soulignant le
rôle du capital social dans la migration et les transferts. Synthétisée par
\textcite{massey1993}, elle établit que les réseaux migratoires, c'est-à-dire les
liens interpersonnels qui relient les migrants installés, les candidats à la
migration et les non-migrants dans les communautés d'origine et de
destination, réduisent considérablement les coûts et les risques du
déplacement~: informations sur les opportunités d'emploi, hébergement à
l'arrivée, aide administrative et financière au voyage. Chaque migration
réussie enrichit le réseau et abaisse le coût des migrations suivantes, ce
qui confère au phénomène un caractère cumulatif et auto-entretenu,
indépendant des différentiels de salaires qui l'ont initié. Au sens de
\textcite{bourdieu1986}, ces réseaux constituent un capital social~: un ensemble
de ressources mobilisables liées à l'appartenance à un groupe, convertibles
en capital économique. Cette lecture a deux implications directes pour les
transferts. D'une part, l'envoi de fonds est le moyen par lequel le migrant
entretient sa position dans le réseau et honore les obligations contractées
envers ceux qui ont facilité son départ, ce qui rejoint les mobiles d'échange
et de dette examinés ci-dessous. D'autre part, la densité du réseau
conditionne les canaux de transfert disponibles et leur coût, les
communautés à forte tradition migratoire, comme celles de la vallée du
fleuve Sénégal, disposant de circuits informels denses et fiables. Cette
théorie éclaire aussi la sélection des ménages bénéficiaires~: recevoir des
transferts suppose d'appartenir à un réseau migratoire, ce qui distingue
systématiquement les ménages récipiendaires des autres et fonde le problème
d'endogénéité traité par la stratégie d'identification de ce mémoire.

\subsubsection{Les motivations de transfert des migrants}

La littérature économique propose plusieurs
cadres pour expliquer pourquoi les migrants transfèrent une partie de leurs
revenus. Ces motivations, concurrentes mais non mutuellement exclusives,
peuvent être regroupées en quatre familles \parencite{rapport2006, azam2016}~:
l'altruisme, l'intérêt personnel et l'échange, la diversification du risque et
l'assurance, et l'investissement collectif. Chaque famille engendre des
prédictions distinctes sur la sensibilité des transferts au revenu du migrant,
au revenu du ménage receveur et à la durée de la migration, prédictions qui
importent pour ce mémoire car le mobile dominant conditionne l'usage des fonds
reçus au bénéfice des enfants.

La première famille est celle de l'altruisme, dont le cadre le plus ancien est
formulé par \textcite{lucas1985} dans leur étude fondatrice sur le Botswana. Ce
modèle postule que le migrant intègre le bien-être de son ménage d'origine dans
sa propre fonction d'utilité~: il tire une satisfaction directe de
l'amélioration du niveau de vie des membres restés au pays, au même titre que
de sa propre consommation. Il en découle deux prédictions. D'abord, les
transferts réagissent positivement au revenu du migrant~: plus celui-ci est
aisé, plus il peut transférer sans réduire sa propre consommation en deçà d'un
seuil acceptable. Ensuite, et c'est la prédiction la plus discriminante, les
transferts réagissent négativement au revenu du ménage receveur~: un ménage
d'origine plus pauvre doit recevoir davantage, puisque l'utilité marginale d'un
franc supplémentaire y est plus élevée. Les transferts jouent alors un rôle de
redistribution intrafamiliale et, par extension, d'assurance implicite contre
la pauvreté du ménage d'origine. Cette sensibilité négative au revenu du
receveur soulève toutefois une difficulté d'identification~: le revenu du
ménage receveur est lui-même en partie déterminé par les transferts déjà
perçus, ce qui complique l'estimation causale du motif altruiste.

La deuxième famille regroupe les mobiles d'intérêt personnel et d'échange, qui
prennent le contre-pied de l'altruisme en faisant du migrant un investisseur
plutôt qu'un donateur désintéressé \parencite{rapport2006}. Dans sa version
patrimoniale, le migrant transfère pour préserver ses droits sur les actifs
familiaux (héritage, terre, logement) ou pour entretenir sa réputation et ses
liens sociaux en vue d'un retour futur~; les transferts dépendent alors surtout
de sa capacité contributive et croissent avec les enjeux de succession. Dans sa
version contractuelle, le migrant et son ménage sont liés par une dette
implicite~: le ménage a financé les coûts initiaux de la migration (scolarité,
voyage, installation) et les transferts en constituent le remboursement, au
même titre qu'un prêt assorti d'un intérêt, ou la rémunération de services
rendus (garde des enfants restés au pays, entretien d'un bien, prise en charge
des parents âgés). Ces mobiles partagent une même logique d'échange
intéressé et une prédiction distinctive~: contrairement à l'altruisme, où les
transferts perdurent tant que le besoin existe, la logique de remboursement
prédit qu'ils décroissent à mesure que la dette est apurée, suivant un profil
en cloche avec la durée de la migration.

La troisième famille rassemble les mobiles de diversification du risque et
d'assurance, au cœur de la nouvelle économie des migrations de travail (NELM)
formulée par \textcite{stark1985}. Elle opère une rupture en substituant à
l'acteur individuel le ménage comme unité de décision~: la migration n'est pas
le choix isolé d'un membre, mais un arrangement collectif visant à diversifier
le portefeuille de revenus du ménage et à le prémunir contre les défaillances
des marchés locaux du crédit et de l'assurance. Les transferts sont alors la
contrepartie de cet arrangement et fournissent une assurance de fait contre les
chocs de revenu agricole ou climatique. Le précautionnement en est le
prolongement direct \parencite{gubert2002}~: les transferts, censés augmenter
en période de choc négatif, permettent au migrant d'agir comme un assureur
résiduel qui complète les marchés formels absents ou incomplets, ou s'y
substitue. Ce cadre explique des faits que l'altruisme pur peine à rendre
compte, notamment la persistance des transferts lorsque le revenu du ménage
s'améliore et le choix stratégique de destinations dont les marchés du travail
sont peu corrélés à celui d'origine. Sa difficulté théorique tient à la
distinction entre un choc transitoire, qui déclenche un transfert ponctuel, et
un choc persistant, qui peut au contraire décourager le transfert.

La dernière famille d'explications, moins explorée dans la littérature mais
particulièrement pertinente pour le Sénégal, envisage les transferts comme un
investissement stratégique et collectif~: elle les analyse non plus comme un
flux bilatéral entre le
migrant et son ménage, mais comme un instrument d'investissement collectif
dans le capital social et physique de la communauté d'origine. Les associations
de migrants, tontines, fédérations villageoises, associations de
ressortissants organisées par région ou par village d'origine, collectent des
fonds auprès de leurs membres établis à l'étranger et les affectent au
financement d'infrastructures locales~: puits et forages, mosquées, écoles,
dispensaires, routes d'accès. Ce canal collectif, analysé par \textcite{azam2016}
pour l'Afrique de l'Ouest, se distingue des mobiles individuels
précédents sur un point essentiel~: il peut produire des effets diffus sur
l'ensemble des ménages d'une communauté, y compris ceux qui ne reçoivent aucun
transfert privé direct, par le biais de l'amélioration des infrastructures et
des services publics locaux. Cette dimension collective est particulièrement
documentée au Sénégal, où les associations de ressortissants originaires de la
vallée du fleuve Sénégal et de la région de Kolda financent traditionnellement
des infrastructures communautaires. Cela suggère que l'impact des transferts sur
le bien-être des enfants pourrait ne pas se limiter aux seuls ménages
bénéficiaires directs identifiés par la variable de traitement de ce mémoire, une limite méthodologique discutée en section~\ref{sec:discussion}.

Ces différents mobiles ne sont pas mutuellement exclusifs et leurs prédictions
se recoupent partiellement \parencite{rapport2006}~: un même transfert peut
relever simultanément de l'altruisme, de l'assurance et du remboursement d'une
dette implicite, leur poids relatif variant selon les caractéristiques du
migrant et de son ménage d'origine. Cette pluralité de motivations, dont
aucune ne s'impose a priori, invite à ne pas privilégier un mobile unique et
justifie de traiter le statut de bénéficiaire comme une variable endogène dont
les déterminants doivent être explicitement modélisés, ce que fait la stratégie
d'identification de ce mémoire.

\subsubsection{Canaux théoriques d'action des transferts sur les dimensions de la pauvreté des enfants}

Ce dernier volet théorique identifie, dimension par dimension, les mécanismes
par lesquels les transferts de migrants peuvent agir sur la pauvreté
multidimensionnelle des enfants. Ces canaux guident la construction des
variables de résultat et l'interprétation des effets empiriques présentés
ensuite.

La théorie des capabilités \parencite{sen1999}, dont l'opérationnalisation en
mesures de pauvreté multidimensionnelle est détaillée dans la section
méthodologique, et les travaux empiriques permettent d'identifier
six canaux principaux par lesquels les transferts de migrants agissent sur
la pauvreté multidimensionnelle des enfants. Ces canaux correspondent, à peu
de choses près, aux dimensions de l'indice N-MODA construit dans ce mémoire,
ce qui permettra d'interpréter les effets estimés par dimension à la lumière
des mécanismes décrits ici.

Les deux premiers canaux relèvent de l'investissement direct en capital
humain. En matière d'éducation, les transferts lèvent la contrainte de
liquidité qui empêche les ménages pauvres d'investir dans la scolarisation
même lorsque ses rendements privés sont positifs~: ils financent les frais de
scolarité, les uniformes et les fournitures, tout en réduisant le recours au
travail des enfants qui concurrence la fréquentation scolaire
\parencite{cox2003, acosta2011}, conformément au modèle altruiste de transfert
\parencite{lucas1985}. Sur le plan sanitaire, l'accroissement du revenu
disponible permet aux ménages bénéficiaires d'accéder aux consultations
médicales, aux médicaments et à la vaccination, dépenses souvent sacrifiées
en premier lorsque le budget se contracte \parencite{hildebrandt2005,
amuedo2007}. Un canal nutritionnel prolonge ce mécanisme sanitaire~: les
transferts modifient directement la structure de la consommation alimentaire,
en augmentant les dépenses en aliments diversifiés et en améliorant le suivi
de la croissance des jeunes enfants \parencite{calero2009}, ce qui touche la
dimension de privation la plus critique pour la petite enfance.

Les trois canaux suivants passent par l'investissement dans l'habitat et son
environnement immédiat. S'agissant de l'accès à l'eau, une partie des fonds
reçus est affectée à des investissements dans le logement, notamment au
raccordement au réseau d'eau ou à l'acquisition d'équipements de stockage
améliorés \parencite{yang2008, gubert2002}. Pour l'assainissement, les
transferts financent la construction ou l'amélioration des toilettes et
latrines, qui donnent accès à des infrastructures sanitaires conformes aux
normes internationales JMP \parencite{azam2016, combes2019}. Enfin, motivés
par un souci d'assurance et de prestige social \parencite{gubert2002,
stark1985}, les ménages bénéficiaires investissent dans la rénovation des
matériaux de construction et l'agrandissement du logement, ce qui réduit le
surpeuplement et améliore la salubrité du cadre de vie des enfants. Ces trois
canaux ont en commun d'exiger des montants transférés suffisamment importants
et réguliers pour financer des investissements indivisibles~: on peut donc
s'attendre à ce qu'ils réagissent moins vite aux transferts que les canaux de
consommation courante, et qu'ils restent conditionnés à l'existence d'une
offre locale (réseau d'adduction, artisans, matériaux) que les fonds privés
ne peuvent créer à eux seuls.

Ces canaux ne sont pas mutuellement exclusifs~: un ménage bénéficiaire peut
simultanément améliorer son alimentation, rénover son logement et financer
la scolarisation de ses enfants. L'intensité relative de chaque canal dépend
du montant, de la régularité des transferts et du contexte socio-économique
du ménage. L'approche multidimensionnelle adoptée dans ce mémoire permet
précisément de capturer l'ensemble de ces effets et d'évaluer leur ampleur relative.

% ── PARTIE II : Revue empirique ────────────────────────────
\subsection{Revue empirique}

Cette revue empirique recense les résultats de la littérature sur
l'impact des transferts de migrants sur le bien-être des enfants. Elle examine
les effets documentés sur chaque grande dimension de la pauvreté
multidimensionnelle (pauvreté monétaire, éducation, santé, logement),
avant de dresser un bilan des travaux portant directement sur la pauvreté
multidimensionnelle et d'identifier les lacunes que ce mémoire cherche à combler.

\subsubsection{Effets empiriques des transferts par dimension de la pauvreté des enfants}

Cette sous-section recense les résultats empiriques en suivant les dimensions
du N-MODA. Elle part de l'effet des transferts sur le revenu et la consommation
des ménages, canal de transmission commun à toutes les dimensions, puis examine
tour à tour l'éducation, la santé, la nutrition, l'habitat (eau, assainissement
et logement) et la protection de l'enfant.

\paragraph{Revenu et consommation des ménages.}
La relation entre transferts de migrants et réduction de la pauvreté monétaire
est l'une des plus documentées dans la littérature du développement.
\textcite{adams2005}, à partir d'un échantillon de 71 pays en développement,
trouvent qu'une augmentation de 10\% de la part des transferts dans le PIB est
associée à une baisse de 3,5\% de la proportion de personnes vivant avec moins
d'un dollar par jour. \textcite{acosta2008}, pour l'Amérique latine, confirment
un effet réducteur de la pauvreté monétaire, bien que l'effet sur les inégalités
soit ambigu en raison de la sélection positive des migrants. Ces résultats
globaux masquent une hétérogénéité considérable~: \textcite{taylor2008}
montrent pour le Mexique que les transferts réduisent la pauvreté dans les zones
à forte émigration mais peuvent accroître les inégalités, et \textcite{azam2016}
concluent, pour l'Afrique, à un effet positif mais modeste, fortement atténué
par les coûts de transaction élevés des transferts formels. Pour l'Afrique de
l'Ouest, \textcite{combes2019} montrent, sur données de panel de l'UEMOA, que
les transferts réduisent la volatilité de la consommation et la pauvreté
transitoire, mais agissent plus faiblement sur la pauvreté chronique. Au
Sénégal, \textcite{fall2010} suggère que les ménages bénéficiaires ont une
consommation par tête supérieure de 15 à 25\% à celle des non-bénéficiaires,
écart partiellement imputable à des caractéristiques non observables. Plus
récemment, \textcite{ndiaye2024} corrigent le biais de sélection sur données
sénégalaises et montrent que la migration et les transferts réduisent la
participation au marché du travail des membres restés au pays tout en favorisant
leur accumulation de capital humain, ce qui souligne la nécessité, retenue dans
ce mémoire, de corriger ce biais avant toute interprétation causale.
\textcite{carrasco2023} estiment enfin, sur les enquêtes longitudinales MAFE et
MESE relatives aux migrants sénégalais en Espagne, des modèles de survie de la
probabilité de transférer~: l'insertion sur le marché du travail du migrant en
est le déterminant le plus robuste et, une fois initiés, les transferts sont
rarement interrompus, ce qui légitime l'hypothèse d'un statut de bénéficiaire
relativement stable entre deux vagues d'enquête rapprochées.

\paragraph{Éducation.}
La littérature sur la scolarisation est abondante mais contrastée.
\textcite{cox2003}, sur le Salvador, montrent qu'un accroissement des transferts
réduit significativement l'abandon scolaire, surtout au secondaire, par la levée
de la contrainte de liquidité. \textcite{acosta2011} confirment ce résultat pour
le même pays et montrent que les transferts réduisent aussi le travail des
enfants, ce qui libère du temps pour l'école. Pour les Philippines,
\textcite{yang2008} exploite les chocs de change comme variation exogène des
transferts et trouve un effet positif sur la fréquentation scolaire,
particulièrement pour les ménages les plus pauvres. En Afrique subsaharienne,
\textcite{gyimah2013} trouvent, pour le Ghana, des taux de scolarisation
primaire plus élevés chez les bénéficiaires, surtout en milieu rural, tandis que
\textcite{azam2016} soulignent que cet effet reste conditionnel à l'existence
d'écoles accessibles et de qualité. Dans le cas sénégalais, la coexistence de
l'école publique et des \textit{daara} complexifie la mesure de la
scolarisation \parencite{unicef2016}. Pour l'Équateur, \textcite{bucheli2018}
nuancent ce tableau~: par un probit bivarié corrigeant l'endogénéité, ils
montrent que l'effet sur la scolarisation secondaire varie fortement selon le
genre, le milieu et la richesse (positif pour les garçons pauvres urbains,
négatif pour les filles rurales, nul pour les ménages aisés), ce qui illustre
l'importance de désagréger les effets, logique reprise dans l'analyse
d'hétérogénéité de ce mémoire.

\paragraph{Santé.}
\textcite{hildebrandt2005} ont été parmi les premiers à documenter, sur données
mexicaines, une mortalité infantile plus faible, un poids à la naissance plus
élevé et une meilleure couverture vaccinale chez les enfants des ménages
bénéficiaires, effets passant par l'accroissement des dépenses de santé.
\textcite{amuedo2007} confirment, pour le Mexique, un effet positif sur les
dépenses de santé préventive, plus marqué pour les ménages ruraux et à faibles
revenus. \textcite{calero2009} trouvent, pour l'Équateur, que les transferts
augmentent les dépenses de santé, mais conditionnellement à l'existence de
services accessibles, ce qui pose la question de la complémentarité entre
transferts privés et offre publique. \textcite{fontenla2022} confirment cet
effet pour les jeunes enfants équatoriens en instrumentant la réception de
transferts, mais montrent qu'il se concentre sur la moitié la plus aisée de la
distribution et sur les garçons. En Afrique de l'Ouest, \textcite{millogo2024}
confirment, pour le Burkina Faso et par un modèle à traitement endogène, un
effet positif sur les indicateurs de santé infantile, ce qui renforce la
plausibilité d'un canal similaire au Sénégal.

\paragraph{Nutrition.}
Sur la dimension nutritionnelle, \textcite{davis2016} exploitent les enquêtes
LSMS guatémaltèques et instrumentent la migration du chef de ménage par la
prévalence migratoire historique~: l'effet des transferts sur les indicateurs
anthropométriques des enfants restés au pays est positif lorsque le père migrant
reste rattaché au ménage nucléaire, mais s'atténue voire s'inverse lorsque son
absence s'accompagne d'une recomposition du ménage. Ce résultat souligne la
nécessité de distinguer l'effet du transfert monétaire de celui de l'absence
parentale qui l'accompagne, distinction que la stratégie PSM-DD de ce mémoire
ne peut lever et qui est discutée comme limite en
section~\ref{sec:discussion}. Dans un registre associant santé et éducation,
\textcite{bouoiyour2016} montrent, pour le Maroc, un effet globalement positif
sur l'accumulation de capital humain des enfants~: les canaux sanitaire,
nutritionnel et éducatif s'articulent souvent au sein d'un même processus
d'investissement familial.

\paragraph{Eau, assainissement et logement.}
Les transferts influencent aussi les conditions d'habitat, dimensions
importantes de la pauvreté multidimensionnelle des enfants. \textcite{yang2008}
montre que les transferts accroissent les dépenses de rénovation du logement des
ménages philippins, et \textcite{combes2019} que la stabilisation de la
consommation en période de choc maintient un accès minimal aux biens essentiels.
Dans le contexte africain, plusieurs travaux montrent que les transferts servent
fréquemment à améliorer le logement (construction, rénovation) et à accéder à
l'eau potable et à l'assainissement \parencite{gubert2002, azam2016}. Ces
investissements restent toutefois conditionnés à l'existence d'une offre locale
que les fonds privés ne créent pas, ce qui limite leur effet sur ces dimensions
lorsque les réseaux d'adduction ou d'assainissement font défaut.

\paragraph{Protection de l'enfant.}
La dimension protection, qui recouvre l'enregistrement des naissances, le travail
des enfants et la séparation d'avec les parents, est la plus directement exposée
aux coûts non monétaires de la migration. \textcite{mansuri2006} trouve, pour le
Pakistan, que la migration du père peut réduire la scolarisation des enfants,
notamment des filles, en l'absence du tuteur principal, effet d'absentéisme
parental d'autant plus pertinent au Sénégal que la migration y est
majoritairement masculine. \textcite{antman2012} signale de même que l'absence
du parent migrant peut nuire à la santé mentale des enfants. La revue de
\textcite{cortes2007} pour l'UNICEF systématise cette distinction entre effet de
la migration et effet des transferts~: à partir des cas dominicain et
salvadorien, elle montre que l'absence du parent perturbe la scolarité et le
suivi de l'enfant, tandis que les transferts reçus exercent un effet
compensateur (une hausse de 1\,\% des transferts réduit l'analphabétisme des
6-14 ans de 5,4\,\% au Mexique selon López, cité par \textcite{cortes2007}).
Cette même synthèse relève, dans plusieurs études de cas, un risque accru de
travail précoce et de vulnérabilité des enfants confiés à des tiers, ce qui
justifie de suivre la protection comme une dimension à part entière du N-MODA.

\subsubsection{Effets ambigus et limites : la mise en garde de la littérature qualitative}

Les canaux d'impact recensés ci-dessus dressent un tableau largement positif
des transferts sur le bien-être des enfants. Cette lecture doit cependant être
nuancée par la synthèse de littérature académique et institutionnelle réalisée
par \textcite{cortes2007} pour la Division des politiques et de la planification de
l'UNICEF, qui met en garde contre une vision trop unilatéralement favorable.
S'appuyant sur des études de cas menées en Équateur, au Mexique, aux Philippines
et en Syrie, cette synthèse souligne d'abord que la migration et les transferts
sont indissociables~: si les transferts améliorent le bien-être matériel de
l'enfant, la migration d'un ou des deux parents qui les génère met simultanément
en péril son développement affectif et psychologique, du fait du manque
d'encadrement parental. Plusieurs études de cas anthropologiques recensées par
\textcite{cortes2007} signalent un risque accru de traite et d'exploitation des
enfants livrés à la garde d'adultes autres que leurs parents, particulièrement
documenté en Albanie et en Moldavie.

La dimension de genre constitue un second point d'attention. \textcite{cortes2007}
montre que la migration féminine, si elle peut accroître l'autonomie économique
des femmes migrantes, expose simultanément les femmes restées au pays et les
migrantes elles-mêmes à un ``triple désavantage'' cumulant inégalités de genre,
de race et de statut migratoire (UNRISD, cité par \textcite{cortes2007}). Dans les
ménages où la mère migre, les pères assument de nouvelles responsabilités
domestiques et éducatives, avec des effets contrastés selon les études sur la
qualité de la prise en charge des enfants. Enfin, cette synthèse relève des
effets économiques et sociaux plus larges susceptibles d'affecter indirectement
le bien-être des enfants~: les transferts peuvent accroître les inégalités de
revenu au sein d'une même communauté, exercer une pression sur les ménages
non-bénéficiaires pour qu'ils envoient à leur tour un membre à l'étranger, et
modifier les habitudes de consommation au point de stigmatiser les enfants des
ménages bénéficiaires dans les communautés rurales aux normes culturelles
établies. La littérature reste également divisée sur l'existence d'une
``culture de dépendance'' aux transferts qui découragerait l'initiative
individuelle des ménages receveurs, un effet que certains auteurs contestent.
Ces mises en garde ne remettent pas en cause l'intérêt d'étudier l'effet des
transferts sur la pauvreté multidimensionnelle des enfants, mais invitent à
une interprétation prudente d'un éventuel effet nul ou modeste dans les
résultats de ce mémoire~: celui-ci pourrait refléter non l'absence d'effet du
transfert monétaire, mais la compensation partielle d'un effet positif du
transfert par les coûts psychosociaux de l'absence parentale qui l'accompagne.

\subsubsection{Synthèse générale de la revue, bilan et lacunes}

Cette synthèse rassemble les enseignements de la revue théorique et de la
revue empirique, avant de recenser les rares travaux portant sur la pauvreté
multidimensionnelle dans son ensemble et de situer la contribution du mémoire.

Un premier enseignement concerne les raisons de la migration et de l'envoi
de fonds, qui commandent la lecture de tout ce qui suit. La théorie
néoclassique explique le départ par les différentiels de rémunération et
l'investissement en capital humain du migrant, tandis que la théorie des
réseaux montre que le capital social des communautés d'origine abaisse le coût
du déplacement et rend la migration cumulative. Ces cadres se prolongent dans
les six motifs de transfert recensés (altruisme, intérêt personnel,
diversification du risque, échange et dette, assurance, investissement
collectif), dont la coexistence a deux conséquences directes. D'une part, le
statut de bénéficiaire n'est jamais aléatoire~: il traduit l'appartenance à un
réseau migratoire et une sélection à la fois positive (ménages plus éduqués,
mieux insérés) et négative (départ du membre le plus apte), ce qui fonde le
problème d'endogénéité au cœur de ce mémoire. D'autre part, le mobile dominant
conditionne l'usage des fonds et donc les dimensions du bien-être de l'enfant
susceptibles d'être affectées.

Un second enseignement porte sur le caractère mitigé des effets lorsqu'on
examine les indicateurs séparément. Sur la pauvreté monétaire, l'effet
réducteur est robuste mais atténué par les coûts de transaction et la
sélection des ménages migrants \parencite{adams2005, azam2016}. Sur
l'éducation, les transferts lèvent la contrainte de liquidité et réduisent le
travail des enfants \parencite{cox2003, acosta2011}, mais l'effet s'inverse
selon le genre et le milieu \parencite{bucheli2018} et l'absence du parent
migrant peut le compenser \parencite{mansuri2006, cortes2007}. Sur la santé et
la nutrition, l'amélioration de l'accès aux soins \parencite{hildebrandt2005,
millogo2024} coexiste avec un effet nul, voire négatif, lorsque l'absence
parentale s'accompagne d'une recomposition du ménage \parencite{davis2016} ou
se concentre sur les ménages déjà aisés \parencite{fontenla2022}. Sur le
logement et les services de base, l'effet dépend d'une offre publique
préexistante que les fonds privés ne créent pas. Il ressort de ce bilan qu'un
effet favorable sur une dimension isolée coexiste fréquemment avec un effet nul
ou défavorable sur une autre, ce qui plaide précisément pour une mesure
\emph{multidimensionnelle} agrégeant l'ensemble de ces canaux plutôt qu'un
indicateur unique.

Trois enseignements méthodologiques en découlent pour ce mémoire.
Premièrement, la quasi-totalité des études crédibles mobilisent une stratégie
de correction de l'endogénéité (variable instrumentale, traitement endogène,
correction de sélection ou appariement)~: la comparaison naïve
bénéficiaires/non-bénéficiaires est disqualifiée. Deuxièmement, les effets sont
systématiquement hétérogènes selon le genre, le milieu, la richesse du ménage
ou le statut migratoire du parent, ce qui justifie l'analyse d'hétérogénéité
conduite au chapitre~\ref{sec:discussion}. Troisièmement, aucune étude recensée
ne combine mesure multidimensionnelle de la pauvreté infantile et
identification sur panel en Afrique de l'Ouest~: c'est la lacune centrale que
ce mémoire comble.

Les travaux qui analysent directement l'impact des transferts sur la pauvreté
multidimensionnelle, et non sur ses composantes prises isolément, restent
rares. \textcite{acosta2008} construisent un indicateur composite pour
l'Amérique latine et trouvent un effet négatif des transferts sur la pauvreté
multidimensionnelle, mais leurs résultats souffrent de l'endogénéité des
transferts. \textcite{mckenzie2012} passent en revue les défis méthodologiques
posés par cette littérature et concluent que la plupart des estimations sont
biaisées à la hausse en raison de la sélection des ménages migrants.


La principale difficulté méthodologique de cette littérature tient au biais
de sélection : les ménages qui reçoivent des transferts diffèrent
systématiquement des ménages non-bénéficiaires, en termes de caractéristiques
observables (niveau d'éducation, localisation, accès au crédit) et inobservables
(réseaux sociaux, propension à migrer, qualité des liens familiaux, aversion
au risque). Une estimation naïve de l'effet des transferts sur la pauvreté
infantile est donc biaisée à la fois par la sélection positive (migrants plus
éduqués) et la sélection négative (familles envoyant le migrant le plus apte).

Plusieurs stratégies d'identification ont été utilisées dans la littérature.
La variable instrumentale est la plus robuste théoriquement : \textcite{yang2008}
exploite les chocs exogènes sur les taux de change, et \textcite{acosta2011}
utilisent la densité historique des réseaux de migrants comme instrument.
Ces stratégies, bien qu'efficaces, requièrent des instruments crédibles qui
ne sont pas facilement disponibles dans le contexte sénégalais.

La régression sur discontinuité a été utilisée par certains travaux
exploitant des seuils réglementaires d'accès aux transferts ou de droit à la
migration \parencite{mckenzie2012}, mais son application reste limitée à des
contextes spécifiques.

L'appariement par score de propension (PSM), introduit par
\textcite{rosenbaum1983} et développé par \textcite{heckman1997}, constitue
une alternative crédible lorsque des instruments exogènes valides ne sont pas
disponibles. Le PSM permet de contrôler les différences observables entre
bénéficiaires et non-bénéficiaires, en construisant un groupe de comparaison
statistiquement comparable au groupe traité sous l'hypothèse d'indépendance
conditionnelle. \textcite{caliendo2008} offrent un guide méthodologique complet
pour sa mise en œuvre et le test de la qualité de l'appariement.

L'estimateur PSM-DD combine ces deux approches : le PSM corrige le
biais de sélection sur les observables, tandis que la double différence (DD)
neutralise les effets fixes inobservables stables dans le temps
\parencite{heckman1998, smith2005}. Cette approche combinée a été utilisée avec
succès dans l'évaluation de politiques sociales en Afrique \parencite{card1994,
angrist2009}, notamment dans les programmes conditionnels de transferts
monétaires, mais reste peu exploitée dans la littérature sur les remittances
privés en Afrique subsaharienne.


La revue de la littérature met en évidence plusieurs lacunes importantes que
ce mémoire entend combler.

\textit{Premièrement}, la pauvreté multidimensionnelle des enfants est rarement
l'objet direct de l'analyse dans les travaux sur les transferts : la plupart des
études se concentrent sur une dimension isolée (scolarisation, mortalité infantile,
dépenses de santé). Or, les enfants subissent des privations simultanées dans
plusieurs dimensions, et l'effet des transferts sur ces privations peut être
très hétérogène.

\textit{Deuxièmement}, très peu de travaux portent sur le Sénégal spécifiquement,
alors que ce pays présente des caractéristiques remarquables : un niveau de
dépendance aux transferts parmi les plus élevés d'Afrique subsaharienne, une
diaspora ancienne et géographiquement diversifiée, et une pauvreté infantile encore
très élevée malgré une croissance économique soutenue.

\textit{Troisièmement}, la combinaison PSM et double différence sur des données de
panel, qui permet de contrôler simultanément les biais de sélection observables
et les effets fixes inobservables, n'a pas, à notre connaissance, été appliquée
à cette question dans le contexte sénégalais. L'existence de deux vagues de
l'EHCVM (2018-2019 et 2021-2022) offre une opportunité inédite pour mettre en
œuvre cette stratégie.

Ce mémoire contribue donc à la littérature sur au moins trois points.
Premièrement, il construit et analyse un indice de pauvreté
multidimensionnelle spécifiquement adapté aux enfants sénégalais, en suivant
la méthodologie N-MODA officielle de l'ANSD et de l'UNICEF. Deuxièmement, il
exploite la dimension panel des données EHCVM I et II, les mêmes ménages
observés avant et après, pour mettre en œuvre un estimateur PSM-DD qui
contrôle simultanément la sélection sur les caractéristiques observables et
les facteurs inobservables stables dans le temps, ce qu'aucune des études
sénégalaises recensées ne fait. Troisièmement, il analyse systématiquement
l'hétérogénéité des effets selon le milieu de résidence, le sexe de l'enfant,
le groupe d'âge et la dimension de privation, répondant ainsi à l'enseignement
transversal de la revue empirique selon lequel l'effet moyen des transferts
masque des effets contrastés selon les sous-populations.
